\section{Related Work} \label{sec:related}

\subsection{Speech and Audio Quality Prediction.}
Automatic speech naturalness assessment has been widely studied through non-intrusive MOS predictors that regress perceptual scores directly from audio. Early neural models such as MOSNet~\cite{lo2019mosnet} demonstrated the feasibility of predicting naturalness without reference signals, while subsequent work improved robustness and coverage across speech synthesis and telephony scenarios, including NISQA~\cite{mittag2021nisqa} and DNSMOS~\cite{reddy2021dnsmos}. More recent approaches, such as UTMOSv2~\cite{baba2024utmosv2}, leverage large-scale self-supervised speech representations to better match human judgments across domains. Beyond speech-specific MOS prediction, audio aesthetics and quality models aim to assess broader perceptual attributes across speech, music, and sound, such as Meta-AES~\cite{tjandra2025meta} and the AudioMOS benchmark~\cite{huang2025audiomos}. Despite strong empirical performance, these models operate as black-box regressors that output scalar scores without interpretable justification, motivating the need for generative reward models that explicitly reason over paralinguistic evidence.

\subsection{Generative Reward Models.}
Reward modeling has recently shifted from discriminative scoring toward \emph{generative} formulations that produce explicit critiques or reasoning before assigning rewards. This paradigm includes LLM-as-a-judge approaches~\cite{li2025generation,li2024llms} and generative verifiers that cast reward modeling as a reasoning task~\cite{zhang2024generativeverifiers}. DeepSeek further introduces inference-time scaling for generalist generative reward models, demonstrating that aggregating multiple independently generated judgments significantly improves robustness and generalization~\cite{liu2025inference}. Related work such as GenRM~\cite{mahan2024generative} and J1~\cite{whitehouse2025j1} shows that generative reward models can be iteratively improved via self-generated rationales and online reinforcement learning. While these approaches are now well-established in text domains, analogous techniques for speech and audio remain under-explored.

In the speech domain, AudioJudge~\cite{manakul2025audiojudge} evaluates the capability of frontier speech LLMs to judge audio quality and paralinguistic attributes through prompt engineering, revealing that most models struggle to provide reliable judgments. The most closely related concurrent works to GSRM are SpeechJudge~\cite{zhang2025speechjudge} and WavReward~\cite{ji2025wavreward}, which train speech LLMs to generate chain-of-thought rationales for pairwise speech comparisons. While effective, both approaches rely on a teacher speech LLM to generate reasoning, which can inherit whatever the teacher model has encoded about naturalness, despite its reliability. In contrast, GSRM decomposes reward modeling into explicit acoustic feature extraction followed by feature-grounded generative reasoning, aligning more closely with how humans assess fine-grained paralinguistic cues.

\subsection{RLHF for Speech and Audio.}
Applying reinforcement learning from human feedback (RLHF) to speech and audio generation is challenging due to the continuous and high-dimensional nature of audio signals, yet recent work has demonstrated promising progress, primarily in text-to-speech (TTS) settings. SpeechJudge~\cite{manakul2025audiojudge} applies direct preference optimization (DPO)~\cite{rafailov2023direct} to post-train TTS models using human preferences. Other studies improve TTS quality using learned or proxy reward signals, including diffusion-based reinforcement learning such as DLPO~\cite{chen2024dlpo} and offline preference optimization for text-to-audio alignment~\cite{liao2024baton}. Additionally, ASR-based rewards have been explored to enhance intelligibility and perceived naturalness of speech model outputs~\cite{liu2025group}. Despite these advances, existing RL-based approaches focus almost exclusively on single-utterance TTS or offline audio generation, and no prior work has investigated leveraging \emph{online} reinforcement learning to improve the naturalness of full-duplex, interactive speech LLMs.
